\section{Lean statements and the path to an iPRA derivation}\label{sec:lean}

\paragraph{Runnable theorems.}
The directory \texttt{lean/} is a standalone standard-library development. The same sources compiled on Lean 4.19.0 and 4.32.0. Both kernel audits cover all 56 public theorems and list only subsets of \texttt{propext} and \texttt{Quot.sound}. In the namespace \texttt{Mertens}, \texttt{sum f n} sums $f$ on $1,\dots,n$, \texttt{divsum f x b} restricts to positive $d\le b$ dividing $x$, and \texttt{delta n} is the indicator of $n=1$. The array \texttt{table n} stores the signed solution of the triangular divisor equations. The value \texttt{mu n} is its last entry, and \texttt{M n} is its prefix sum. The zero arguments of both functions evaluate to zero. This choice makes the cancellation lemma an immediate invariant of table extension. The module \texttt{PrimeStep.lean} also constructs the list of prime factors by bounded trial division and proves directly that its repeated-factor and parity rule gives this function.

\begin{center}\small
\begin{tabularx}{\linewidth}{@{}p{.30\linewidth}X@{}}
\toprule
Proved for all inputs & Lean declarations\\ \midrule
Prime factors, signs and interval bounds & \texttt{factorization}, \texttt{factorMu\_correct}, \texttt{mu\_range}, \texttt{M\_interval}, \texttt{M\_bound\_after\_two}\\
Divisor characterization and first difference & \texttt{mu\_divsum}, \texttt{mu\_unique}, \texttt{M\_step}\\
Hyperbola reindexing and floor recurrence & \texttt{weighted\_floor\_sum}, \texttt{summatory\_swap}, \texttt{quotient\_sum}, \texttt{quotient\_recurrence}\\
Ordered products and tail vanishing & \texttt{hyperboloid}, \texttt{no\_tuples}\\
Totient and summatory convolution & \texttt{phi\_convolution}, \texttt{summatory\_convolution}, \texttt{summatory\_totient}\\
Shared table and certified finite lookup & \texttt{fastM\_correct}, \texttt{computes143}, \texttt{lookup143} (the lookup has the stated finite range)\\
\bottomrule
\end{tabularx}
\end{center}

Run the commands in \texttt{lean/README.md} from that directory. The audit lists each theorem's kernel axiom dependencies. These are ordinary Lean theorems about numbers and finite sums. For a deep embedding the goal is a term of type \texttt{Prf PRAAx [] formula} with a concrete compiled \texttt{formula}. ProveIt's arithmetic library and bridge distinguish this judgment from a semantic equality over \texttt{Nat} or \texttt{Int} \cite{repo}. No assumption asserting one of the analytic results is required for the executable finite layer.

\paragraph{Finite modules to build next.}
For each of the following instructions use the numerical encodings in Section~\ref{sec:foundation}. Fix the largest input index before applying induction to a table valid at several smaller arguments.
\begin{enumerate}
\item Encode \texttt{mu}, \texttt{M}, a product, a bounded sum, and the nonzero-divisor indicator as PR programs. The table invariant states at once that every old entry agrees with its prefix and every completed row of divisor tests has the prescribed sum. Compile that conjunction to a Boolean fold. Prove the two sum-reordering rules using a rectangle clipped by $ad\le N$. Instantiate the weight by $1$, the identity function, and the bounded search $v_p$. This gives the floor, $\varphi$ and Chebyshev endpoints. The tuple-count array has width $N+1$ and height $L_N+1$. Its zero tail is proved within the same finite rectangle.
\item Define the permutation sign by a double loop counting inversions in a base-$(N+1)$ digit list and the determinant by the bounded list sum. Prove the row-swap pairing of Theorem~\ref{thm:redheffer} on that list, then specialize to row 1. For Farey define the reduced-pair indicator and cyclic coefficient convolution on a fixed length $N!$. Implement the polygon identity as coefficient rotation, followed by the monic polynomial divisions and extended gcd in Theorem~\ref{thm:farey}. Every vector equality is a bounded conjunction of signed-coordinate equalities.
\item In the sieve maintain the prime list, composite bits, coefficients, and a prefix up to the current index. In the isolated algorithm the invariant is equality to the defining $M$ at every sieved index and every completed large quotient. The termination measures are the remaining indices and the next positive-length block. Count completed blocks separately from arithmetic correctness. For the squarefree count use the Boolean square-divisor scan and then the indicator and telescoping bounds of Section~\ref{sec:squarefree}.
\end{enumerate}

\paragraph{Compiling the analytic statements.}
Use complex rectangular rational enclosures, with rational lower bounds for denominators. Addition and multiplication enlarge error by the triangle inequality. If $|z|\ge a>0$, evaluating $1/z$ at a centre within $e\le a/2$ incurs at most $2e/a^2$. For an input with bound $B$ on its magnitude this also bounds the required input precision of a finite product by forward propagation. A direct implementation can instead propagate requested radii backwards. No equality test on a complex name is used. For each fixed computation, choose a bound $D\ge2$ for the local Lipschitz constants and for the number of inputs to each operation. At a graph depth $m$, per-node error less than $2^{-j-3}(2D^2)^{-m}$ suffices for final error $2^{-j}$ by induction on layers. All bounds used to choose $D,m$ must be computed from the stated size, tail and separation data.

\begin{enumerate}
\item Define \texttt{powPos} from the $\log$ and exponential routines, and a truncated Dirichlet sum and prime product from the coefficient table. On a fixed input with $\Re s\ge1+1/q$, the geometric majorant gives the truncation at any bit demand. Prove the rectangle product bound, then the cellwise finite Mellin identity. Express both the corrected Perron kernel and its side bounds as a loop over rational quadrature meshes. Split the integer endpoint exactly at the Boolean test $a=bn$.
\item For the squarefree constant, the finite floor estimate is rational. Identify the reciprocal at $2$ by the product lemma. Independently, one can substitute ProveIt's rational Basel estimate for the trigonometric calculation. For $n\ge 2^{2j+7}$ the packet gives $|6\sum_{\ell\le n}\ell^{-2}-P_j^2|\le2^{-j}$ for its fast rational name $P_j$ of $\pi$ \cite{basel}. Both paths use a denominator separated from zero.
\item For the logarithmic saving, implement the cut-off continuation formula with one cell per integer, rational height $T$ and a certified lower bound on the regularized function. Use an explicit grid on the rectangle, second derivative bounds and cancellation of its interior edges. The output of the final comparison can be taken entirely integral, e.g.\ $|M(N)|L^h\le K_h(2N+1)$ with bit length $L$ as in Section~\ref{sec:logrates}. The proof there gives this deliberately weakened integer comparison directly. For the alternative PNT route instantiate the named PR function $F$ and construct the target $[N\ge F(k)\Rightarrow k|M(N)|\le N]$.
\item For the growth criterion, retain the numerical witness $C$, the gap $\Delta$, a bound on $|s|$ and a bit demand. Use the geometric coefficient tail at $z_0=s+2$, the finite Abel coefficient equations, and the convolution proof. To bound the tail of the product past degree $J$, each discarded pair has an index exceeding $\lfloor J/2\rfloor$. Bound each factor's full majorant times the other's half-degree tail. Local residue calculations add a Taylor-coefficient function and the given derivative-apartness certificate.
\item Express compact zero exclusion by a closed first-order sentence with parameters $m,T,b$, an enumeration of rational points in the rectangle, and a bit demand. Use the disk quadratures of Section~\ref{sec:compact-rh}, an enclosing rectangle's witness $b$, the chosen coefficient cutoff $J$, and an explicit largest Fourier frequency. Compute the final growth constant before specializing the bound on compact sets. In the lower-bound argument of Section~\ref{sec:omega} encode the search on $[A,B]$ as a Boolean fold with a first-hit index. If it returns no hit, use its finite invariant to bound the clipped table, instantiate the certified zero name, and apply the two rational Taylor estimates with the stated strict contradiction margin. This proves that the first-hit output satisfies the integer test.
\end{enumerate}

For an object derivation, replace each rational test by the equivalent natural-number comparison after cross multiplication. Prove the one-step rounding and truncation inequalities first. Collect all local obligations up to the fixed precision and cutoff as a bounded array of Booleans, and apply quantifier-free induction to its prefix. A branch concerning a true-real boundary is implemented by rational subdivision with the overlapping certified ranges given in the proof. Quantify the input numerals and bit demand only after constructing that derivation. This supplies the concrete induction predicates needed to move from a semantic Lean file towards ProveIt's deeply embedded logic.
